\documentclass[11pt]{report}
\input{../pri-end/T0_preamble_standalone_A4_De}
% T0_preamble_patches.tex — LOKALER PLATZHALTER
% Im Repo durch die offizielle Version ersetzen.


\title{Dok.~309 -- Das Skalenanker-Problem:\\
$\Lambda$CDM und FFGFT im Vergleich}
\author{Johann Pascher}
\date{Juli 2026}

\begin{document}
\maketitle
\tableofcontents

% ============================================================
\chapter*{Dok.~309 -- Das Skalenanker-Problem}
\addcontentsline{toc}{chapter}{Dok.~309 -- Das Skalenanker-Problem}
% ============================================================

% ============================================================
\section*{Worum es geht}
% ============================================================

Beide kosmologischen Theorien — $\Lambda$CDM und FFGFT — müssen irgendwo
eine Zahl einführen, die den kosmischen Maßstab an das SI-System ankoppelt.
Keine der beiden leitet diesen Anker aus erster Prinzipien her.

Das ist kein Zufall — es ist ein strukturelles Problem jeder physikalischen
Theorie, die Dimensionen hat. Aber \textbf{Art und Qualität} des Ankers
unterscheiden sich fundamental. Und es gibt eine Asymmetrie, die bisher
in keinem Dokument explizit benannt wurde: FFGFT erbt durch seinen
Skalenfit die Systematik der Expansionslesart.

Das Dokument ist kein Teil der A-Serie.
Bezüge: Dok.~190 (P20, P33), Dok.~267 (Entartung), Dok.~279 ($H_0$-Kette),
Dok.~308 (kosmischer Sektor).

% ============================================================
\section*{Konventionen}
% ============================================================

\noindent\textbf{Zeichenerklärung.}

\medskip
\begin{tabular}{@{}lp{11cm}}
\textbf{[K]}       & \emph{Kernableitung} — mathematisch oder numerisch
                     aus den Grundsetzungen hergeleitet; reproduzierbar. \\[4pt]
\textbf{[S]}       & \emph{Strukturaussage} — begründet, aber noch nicht
                     vollständig bewiesen. \\[4pt]
\textbf{[SETZUNG]} & \emph{Setzung} — deklarierte Annahme, nicht weiter
                     hergeleitet. \\
\end{tabular}

\medskip
\noindent\textbf{Zentrale Größen.}

\medskip
\begin{tabular}{@{}lp{10.5cm}}
$\xi = 4/30000$
  & Einziger Fundamentalparameter der FFGFT [SETZUNG, P33].
    Form ($4/3$, $3$, $2^2$) vorwärts hergeleitet [K];
    Magnitude ($5^4$) intuitiv/empirisch [SETZUNG]. \\[4pt]
$\bar\lambda_e = \hbar/(m_e c)$
  & Reduzierte Compton-Wellenlänge des Elektrons;
    $= 3{,}8616\times10^{-13}$~m. \\[4pt]
$H_0$
  & Kosmologischer Parameter; in FFGFT:
    $H_0 = (\pi/2)\,c\,\xi^{10}/\bar\lambda_e = 66{,}82$~km/s/Mpc
    [K], Exponent~10 [SETZUNG, P20]. \\[4pt]
$R_H = c/H_0$
  & Geometrische Maximalskala der kompakten Geometrie —
    kein Expansionshorizont (Dok.~182). \\[4pt]
$\nu_\text{Cs}$
  & Caesium-Hyperfeinübergang; $= 9{,}192\,631\,770\times10^9$~Hz
    (exakt, Definition); konventioneller SI-Anker. \\[4pt]
$\nu_e = m_e c^2/h$
  & Elektronenfrequenz; $= 1{,}235\times10^{20}$~Hz;
    lesartunabhängige Bezugsfrequenz. \\[4pt]
$n$
  & Ganzzahliger Exponent in $H_0 \propto \xi^n/\bar\lambda_e$;
    $n = 10$ ist Setzung (P20, Dok.~190). \\[4pt]
$\Omega_\Lambda$
  & Dunkle-Energie-Dichte in $\Lambda$CDM;
    ohne lesartunabhängigen Referenten in der statischen Lesart. \\
\end{tabular}

% ============================================================
\chapter{Das SI-System — ein Netz ohne ersten Anker}
% ============================================================

\section{Alle Einheiten auf eine zurückführbar}

Seit 2019 sind alle SI-Basiseinheiten über fixierte Naturkonstanten
definiert. Die scheinbar „vier unabhängigen Messwerte" $c$, $h$, $e$,
$k_B$ sind physikalisch aber keine unabhängigen Größen — sie sind alle
auf einen einzigen historisch-konventionellen Anker zurückführbar:
die Frequenz des Caesium-Hyperfeinübergangs $\nu_\text{Cs}$.

\begin{center}
\begin{tabular}{lll}
\toprule
Einheit & Rückführung & Historischer Ursprung \\
\midrule
Meter   & $c/\nu_\text{Cs}$ (Zahlenfaktor) & Erdmeridian \\
Kilogramm & $h\,\nu_\text{Cs}^2/c^2$ & Urkilogramm \\
Ampere  & $e\,\nu_\text{Cs}$ & Kraftdefinition \\
Kelvin  & $h\,\nu_\text{Cs}/k_B$ & Thermometrie \\
\bottomrule
\end{tabular}
\end{center}

\section{Es gibt keinen ersten Anker}

$\nu_\text{Cs}$ wurde 1967 gewählt weil es technisch stabil
realisierbar war — nicht weil es physikalisch ausgezeichnet ist.
Der Rb-Übergang, die Elektronenfrequenz oder hundert andere Systeme
hätten ebenso funktioniert. Alle heutigen SI-Zahlenwerte wären dann
andere — alle dimensionslosen Verhältnisse ($m_e/m_p$ usw.) blieben
identisch.

\textbf{Was wirklich festgelegt ist:} nur dimensionslose Verhältnisse
wie $m_e/m_p \approx 1/1836$. Diese sind einheitenunabhängig.

\textbf{Konsequenz:} Jede Theorie, die kosmologische Absolutwerte in SI
voraussagt, gibt implizit eine Relation zwischen zwei Frequenzen an —
und muss dafür irgendwo eine Zahl setzen. Die FFGFT-Aussage

\[
  H_0 = \frac{\pi}{2}\,\frac{c\,\xi^{10}}{\bar\lambda_e}
\]

ist äquivalent zur dimensionslosen Aussage:

\[
  \frac{H_0}{\nu_e} = \frac{\pi}{2}\,\xi^{10} = 2{,}79\times10^{-39}
  \qquad\text{mit}\quad \nu_e = \frac{m_e c^2}{h}.
\]

Der SI-Bezug steckt vollständig in der Wahl von $\bar\lambda_e$ als
Bezugsgröße. Ohne diese Wahl ist $H_0/\nu_e$ eine reine Zahl.

% ============================================================
\chapter{$\Lambda$CDM und sein Skalenanker}
% ============================================================

\section{Sechs freie Parameter, alle empirisch}

$\Lambda$CDM beschreibt die Kosmologie mit sechs freien Parametern
(Planck 2018): $H_0$, $\Omega_b h^2$, $\Omega_\text{DM} h^2$, $n_s$,
$\sigma_8$, $\tau$. Alle aus Beobachtung — keiner aus Theorie abgeleitet.

\section{$H_0$: zirkulär gemessen}

$H_0$ wird über eine vierstufige Entfernungsleiter bestimmt:
Parallaxe → Cepheiden → SNe~Ia → Hubble-Gesetz. Jede Stufe bringt
eigene Kalibrierungsannahmen. Das \textbf{Hubble-Spannungsproblem}
(Planck: $67{,}4$, SH0ES: $73{,}0$~km/s/Mpc, $5\sigma$-Spannung) ist
die direkte Folge dieser Zirkularität.

\section{Das Feinabstimmungsproblem}

$\Lambda$ entspricht einer Vakuumenergiedichte
$\rho_\Lambda \approx 5{,}96\times10^{-27}$~kg/m$^3$.
Die QFT sagt $\rho_\text{QFT} \sim 5\times10^{96}$~kg/m$^3$ vorher.
Das Verhältnis beträgt $10^{123}$. $\Lambda$CDM hat keine Erklärung
dafür — es setzt $\Lambda$ auf den Beobachtungswert ohne
theoretischen Ursprung.

% ============================================================
\chapter{FFGFT und ihr Skalenanker}
% ============================================================

\section{Struktur von $\xi$ [K] und Magnitude [SETZUNG]}

$\xi = 4/30000 = 1/(3\cdot2^2\cdot5^4)$.
Laut P33 (Dok.~190):

\begin{itemize}
  \item \textbf{Form} ($4/3$, Faktor~3, Faktor~$2^2$): vorwärts hergeleitet [K]
  \item \textbf{Magnitude} ($5^4 = 625$): intuitiv/empirisch begründet [SETZUNG]
\end{itemize}

\section{Der eigentliche Fit: der Exponent 10}

$\xi$ selbst kommt aus der Geometrie. Der \textbf{Exponent 10} in
$H_0 = (\pi/2)\,c\,\xi^{10}/\bar\lambda_e$ wurde so gewählt, dass
$H_0 \approx 67$~km/s/Mpc herauskommt. Da $\xi \sim 1{,}3\times10^{-4}$,
springt $H_0$ zwischen benachbarten Exponenten um Faktor $\sim$7500:

\begin{center}
\begin{tabular}{clc}
\toprule
$n$ & $H_0$ & Bewertung \\
\midrule
9  & $5{,}0\times10^5$~km/s/Mpc & unphysikalisch \\
\textbf{10} & $\mathbf{66{,}8}$~\textbf{km/s/Mpc} & plausibel \\
11 & $8{,}9\times10^{-3}$~km/s/Mpc & unphysikalisch \\
\bottomrule
\end{tabular}
\end{center}

Es gibt kein kontinuierliches Anpassen — $n=10$ ist die einzige
ganzzahlige Wahl die einen physikalisch sinnvollen Wert ergibt.
Die Begründung bleibt: $n=10$ gibt einen mit $\Lambda$CDM
kompatiblen Wert (P20, Dok.~190).

\section{Die Vererbung: FFGFT erbt die $\Lambda$CDM-Systematik [K]}

$67{,}4$~km/s/Mpc ist ein $\Lambda$CDM-Messwert — bestimmt unter
der Annahme der Expansionslesart. Wenn $n=10$ an diesen Wert
kalibriert wurde, erbt FFGFT die Systematik dieser Lesart.

Der „natürliche" $H_0$ aus der FFGFT-Rotverschiebungsformel
$z = e^{\xi x/\bar\lambda_e} - 1$ wäre:
\[
  H_0^\text{natur} = \frac{c\,\xi}{\bar\lambda_e}
  \approx 3\times10^{33}~\text{km/s/Mpc}
  \qquad\text{(unphysikalisch)}
\]

Um $\sim$67~km/s/Mpc zu erhalten braucht man $\xi^{10}$ statt $\xi^1$ —
und die Rechtfertigung dieses Exponenten ist letztlich die
$\Lambda$CDM-Kompatibilität.

Das ist keine alleinige Schwäche der FFGFT — es ist die unvermeidliche
Konsequenz der Entartung (Dok.~267): solange beide Lesarten empirisch
ununterscheidbar sind, kalibriert jede Theorie ihre Parameter an
denselben Daten.

Der Ausweg wäre eine lesartunabhängige Herleitung des Exponenten~10
aus der T$^4$-Geometrie vorwärts — bisher nicht geleistet (P20, Dok.~190).

\section{Querverankerung: was FFGFT trotzdem leistet}

Dass $\xi^{10}$ nicht nur $H_0$ sondern auch $a_0$ (MOND-Skala,
Dok.~308) reproduziert — das wäre bei einem reinen $H_0$-Fit nicht
erwartet. Das ist die Querverankerung.

\begin{center}
\begin{tabular}{llll}
\toprule
Sektor & Größe & $\xi$-Abhängigkeit & Status \\
\midrule
Kosmologie & $H_0$ & $\xi^{10}$ & [K] \\
Trägheitsregime & $a_0$ & $\xi^{10}$ & [K]/[S] \\
Teilchenphysik & Leptonmassenleiter & $\xi$-Potenzen & [K] \\
Quantenmechanik & $K_\text{frak} = 1-100\xi$ & $\xi$ & [K] \\
Geometrie & $k^* = \ln\varphi/\xi$ & $\xi$ & [S] \\
\bottomrule
\end{tabular}
\end{center}

$\Omega_\Lambda$ in $\Lambda$CDM hat \textbf{null} solcher Querverankerungen.

% ============================================================
\chapter{Vergleich}
% ============================================================

\begin{center}
\begin{tabular}{lll}
\toprule
Eigenschaft & $\Lambda$CDM & FFGFT \\
\midrule
Freie Skalenanker & 6–7, empirisch & 1 ($\xi$, P33) \\
Theoretischer Ursprung & keiner & keiner \\
Querverankerung & nein & ja (5+ Sektoren) \\
Zirkularität & ja ($H_0$-Leiter) & nein \\
Vererbung $\Lambda$CDM-Systematik & inhärent & via Exponent 10 \\
Feinabstimmung & $10^{123}$ ($\Lambda$) & entfällt strukturell \\
Hubble-Spannung & $5\sigma$, offen & strukturell absent \\
\bottomrule
\end{tabular}
\end{center}

\medskip
FFGFT tauscht sechs isolierte Parameter gegen einen querverankerten —
kein Beweis, aber eine Reduktion mit Richtung. Die Vererbung des
$\Lambda$CDM-Skalenwertes über den Exponenten~10 bleibt die
zentrale offene Schwachstelle.

% ============================================================
\chapter{Offene Punkte (Dok.~190)}
% ============================================================

\begin{itemize}
  \item \textbf{P20:} Form von $H_0$ als $\xi$-Verhältnis fest;
        Exponent~10 ist externe Kalibrierung aus $\Lambda$CDM,
        nicht hergeleitet.
  \item \textbf{P33:} Magnitude von $\xi$ (Faktor $5^4$) intuitiv/
        empirisch begründet, kein Vorwärtsbeweis. Struktur-Parallele
        zu P20.
  \item \textbf{Vererbung (neu, Dok.~309):} Exponent~10 an
        $H_0^\Lambda$CDM kalibriert. Ausweg: lesartunabhängige
        Herleitung aus T$^4$-Geometrie — bisher nicht geleistet.
\end{itemize}

\end{document}
